\chapter{Học để làm gì khi đời không dễ}
\markboth{Học để làm gì khi đời không dễ}{Học để làm gì khi đời không dễ}

\section{Ra đời còn chưa chắc kiếm được việc, học dữ vậy để làm gì?}
\begin{truyen}{Ra đời còn chưa chắc kiếm được việc, học dữ vậy để làm gì?}{Classroom Talk}
\chuhoa{M}{ột bạn chống cằm than: ``Nghe nói tốt nghiệp xong vẫn thất nghiệp đầy. Vậy học cắm đầu giờ để làm gì hả thầy?''}

Thầy đáp: ``Không ai dám hứa học xong là đời dễ. Nhưng không học thì lúc đời khó, em còn ít đồ nghề hơn để tự cứu mình.''

Bạn ấy im một nhịp rồi nói nhỏ: ``Vậy học không phải vé thắng, mà là đỡ tay trắng hả thầy?''

Thầy gật đầu: ``Đúng. Học không cứu em khỏi mọi khổ. Nhưng nó làm em bớt bất lực.''

\textit{(Study is not a guaranteed win. It is a way to be less helpless when life becomes difficult.)}
\end{truyen}

\begin{baihoc}
Không phải cứ học là chắc thắng. Nhưng không chịu học thì thua thường đến sớm hơn.
\textit{(Study does not guarantee success, but refusing to learn often brings failure sooner.)}
\end{baihoc}

\begin{ghinhoanh}
\textbf{Short English to remember:}

Study is not a magic ticket.

It is a tool against helplessness.
\end{ghinhoanh}

\begin{tuhocnhanh}
1. Em đang học để thắng, hay học để đỡ bất lực? \textit{(Do you study to win, or to be less helpless?)}

2. Khi gặp chuyện khó, em có đủ kỹ năng để tự xoay chưa? \textit{(When life gets hard, do you have enough skills to handle it?)}
\end{tuhocnhanh}
\ngancach

\section{Nhà em nói kiếm tiền mới quan trọng, học sau cũng được}
\begin{truyen}{Nhà em nói kiếm tiền mới quan trọng, học sau cũng được}{Classroom Talk}
\chuhoa{C}{ó bạn nói rất thật: ``Nhà em lo tiền từng ngày. Nên nghe ai nói học cho cao siêu em thấy xa lắm. Kiếm tiền mới là việc gấp.''}

Thầy không phản bác ngay. Thầy hỏi: ``Vậy em muốn cả đời chỉ đủ sống từng ngày, hay sau này đỡ bị đồng tiền dắt đi?''

Bạn ấy ngẫm rồi nói: ``Chắc ai nghèo cũng nghĩ tiền gấp trước đã.''

Thầy đáp: ``Đúng. Nhưng nếu chỉ lo cái gấp mà bỏ cái gốc, vài năm nữa em vẫn chạy vòng đó thôi.''

\textit{(Money can be urgent, especially in hardship. But neglecting learning may keep a person trapped in the same cycle.)}
\end{truyen}

\begin{baihoc}
Người khó khăn càng cần học thực chất, vì đó là cách chậm mà bền để không nghèo mãi.
\textit{(Those in hardship need real learning even more, because it is the slower but steadier way out.)}
\end{baihoc}

\begin{ghinhoanh}
\textbf{Short English to remember:}

Urgent money matters.

But skills decide how long money stays.
\end{ghinhoanh}

\begin{tuhocnhanh}
1. Em đang giải quyết cái gấp hay đang xây cái bền? \textit{(Are you solving the urgent, or building the lasting?)}

2. Em muốn học gì để sau này kiếm tiền đỡ bấp bênh hơn? \textit{(What do you want to learn so your future income is less unstable?)}
\end{tuhocnhanh}
\ngancach

\section{Em học hoài vẫn đuối, vậy cố để làm gì?}
\begin{truyen}{Em học hoài vẫn đuối, vậy cố để làm gì?}{Classroom Talk}
\chuhoa{M}{ột bạn bực bội: ``Bạn khác học một hiểu mười. Em học năm hiểu một. Cứ kiểu này thì cố để làm gì?''}

Thầy hỏi lại: ``Em muốn bằng người khác ngay, hay muốn em của tháng sau đỡ tệ hơn em của tháng này?''

Bạn ấy nhăn mặt: ``Nghe vậy thì đỡ cay hơn, nhưng vẫn mệt.''

Thầy nói: ``Mệt là thật. Nhưng mệt vì đang đi lên khác với mệt vì đứng yên rồi tự chê mình.''

\textit{(Progress is not always fast. But there is a difference between being tired from growth and being stuck in self-pity.)}
\end{truyen}

\begin{baihoc}
Cố gắng không phải để chứng minh mình hơn ai. Nhiều khi chỉ là để mai đỡ dở hơn hôm nay.
\textit{(Effort is not always about beating others. Sometimes it is just about being less weak tomorrow.)}
\end{baihoc}

\begin{ghinhoanh}
\textbf{Short English to remember:}

Slow progress is still progress.

Do not compare speed before checking direction.
\end{ghinhoanh}

\begin{tuhocnhanh}
1. Em đang thua người khác hay đang bỏ cuộc với chính mình? \textit{(Are you losing to others, or quitting on yourself?)}

2. Một điều nhỏ nào chứng minh em đã khá hơn trước? \textit{(What small thing proves you are better than before?)}
\end{tuhocnhanh}
\ngancach

\section{Em học chỉ để ba mẹ đỡ chửi, vậy có sao không?}
\begin{truyen}{Em học chỉ để ba mẹ đỡ chửi, vậy có sao không?}{Classroom Talk}
\chuhoa{C}{ả lớp cười khi một bạn nói huỵch toẹt: ``Nói thật nha thầy, nhiều bữa em học không phải vì thích. Em học để khỏi bị la, khỏi bị so với con nhà người ta.''}

Thầy đáp: ``Lý do đó không đẹp, nhưng ít ra còn kéo em ngồi vào bàn. Vấn đề là em có chịu đi tiếp xa hơn lý do đó không.''

Bạn hỏi: ``Nghĩa là sao?''

Thầy nói: ``Ban đầu em học vì sợ. Sau đó em phải dần học vì tương lai của chính em. Nếu đứng yên ở nỗi sợ, em sẽ học trong bực bội suốt.''

\textit{(Fear may start the habit, but maturity means moving from outside pressure to inner purpose.)}
\end{truyen}

\begin{baihoc}
Động lực xấu vẫn tốt hơn không có động lực. Nhưng nếu không nâng nó lên, em sẽ rất nhanh kiệt sức.
\textit{(A weak motivation is still better than none, but if it never grows, burnout comes quickly.)}
\end{baihoc}

\begin{ghinhoanh}
\textbf{Short English to remember:}

Fear can start a habit.

Purpose is what keeps it alive.
\end{ghinhoanh}

\begin{tuhocnhanh}
1. Lý do học hiện tại của em là gì? \textit{(What is your current reason for studying?)}

2. Em muốn đổi từ học vì sợ sang học vì điều gì? \textit{(What do you want to study for besides fear?)}
\end{tuhocnhanh}
\ngancach